\chapter{Approvals}

Most fund losses are interaction mistakes, not stolen
seeds. You signed something that let a contract pull.
The wallet was fine. The spend was not.

\section{Exact amounts}

\texttt{approve()} gives a spender permission to move a
token. Most dApps ask for unlimited. Give the amount
this transaction needs. Revoke the rest later with
Revoke.cash or the explorer's approval checker. Do that
on a schedule, not after a scare.

\section{Permit is not a login}

EIP-2612 \texttt{permit} is an off-chain approval. No
on-chain tx appears until someone else submits it. A
fake ``sign in'' prompt with \texttt{spender},
\texttt{value}, \texttt{nonce}, and \texttt{deadline} is
a permit. You are not logging in. You are loading the
gun.

\section{The usual traps}

\textbf{Address poisoning.} A tiny transfer from an
address that looks like yours, or like a recipient you
use. Scam tokens in the history. Copy from an address
book you maintain, not from ``recent.''

\textbf{Clipboard malware.} It swaps the address you
copied. Read the destination on the hardware screen,
character by character, before you confirm.

\textbf{Fake airdrops.} Unknown tokens in the wallet.
Claiming or swapping them is how
\texttt{setApprovalForAll} gets signed. Ignore them.

\textbf{Ice phishing.} You sign a real \texttt{approve}
to a spender you were talked into. The chain did what
you asked. The social layer was the exploit.

\section{Trades}

Slippage of 5--10\% on a liquid pair is how you get
sandwiched. Tiny slippage is how you burn gas on a
fail. The source's 0.5--1\% band is an example for
liquid pairs, not a law.

Private relay routes (Flashbots Protect, MEV Blocker)
hide the tx from public searchers. Set a deadline so a
stuck tx cannot execute later at a worse price. If an
aggregator hops through a token you did not pick, stop.

\section{Before you click}

Verify the contract. Simulate. Check the approval
amount. Read the typed data. If you do not understand
the message, do not sign it.
